\chapter{Einleitung}
\label{chap:Einleitung}
\section{Relevanz}
In den letzten Jahren hat die Sicherheit im digitalen Raum stark an Bedeutung gewonnen. Firmen werden ständig angegriffen, und man sucht nach neuen Methoden, um private Daten zu schützen. Die Kryptographie wird laufend erweitert und verbessert, gleichzeitig entstehen aber auch neue Methoden, sie zu brechen. Eine grosse Veränderung bringt zurzeit die KI.\\
"Experten warnen schon lange vor Cyberattacken mit KI-Software. [\dots] Die Software mit dem Namen Mythos entdeckte über Jahrzehnte unentdeckt gebliebene Sicherheitslücken bei viel genutzten Programmen und Onlinediensten."\todo{Quelle für das Mythos-Zitat in Zotero ergänzen und hier mit \texttt{\textbackslash citep} einfügen}\\
Die KI hilft dabei, Sicherheitslücken zu schliessen, kann aber ebenso dazu benutzt werden, sogenannte Exploits auszunutzen. Auch die Optimierung des Netzwerks wird immer wichtiger. Heutige Geräte nutzen bereits Wi-Fi 7, und Wi-Fi 8 ist in Entwicklung und könnte in naher Zukunft verfügbar sein.\todo{Quelle für Wi-Fi 7 und Wi-Fi 8} Die Technikbranche wächst stetig, und es werden immer komplexere Systeme entwickelt, etwa in der KI-Forschung, bei Quantenchips und Quantenkryptographie, Natrium-Ionen-Batterien und Bio-Printing.\todo{Quelle für die genannten Technologien} Ich vermute, dass Cybersecurity und Netzwerkanalyse in Zukunft eine immer grössere Rolle spielen werden, auch bei der Frage, wer für Straftaten mit KI verantwortlich ist. Ein Beispiel ist der KI-gesteuerte Cyberangriff auf Hugging Face, für den OpenAI die Verantwortung übernommen hat.\\
"Zunächst brach die Software aus der Testumgebung aus, [\dots] und zwar über eine bisher unentdeckt gebliebene Schwachstelle. Die Modelle kamen den Angaben zufolge eigenständig zu dem Schluss, dass es bei Hugging Face nützliche Daten und Lösungen für die ExploitGym-Aufgaben geben könnte.
Die Software habe sich daraufhin Zugang zu »geheimen Informationen« auf der Plattform verschafft, [\dots]"\citep{sprothenKuenstlicheIntelligenzOpenAI2026}\\
\section{Fragestellung}
Ich wollte herausfinden, wie sich Daten sicher von einem Gerät zu einem anderen verschicken lassen. Die Frage hat zwei Teile. Der erste ist der effiziente Versand von Datenpaketen zwischen zwei Teilnehmern, der zweite die Verschlüsselung, damit nur Teilnehmer mit dem gleichen Schlüssel die Nachrichten lesen können.
Das Ziel dieser Arbeit ist es, mithilfe einer Client-Server-Architektur eine Kommunikation in einem lokalen Netzwerk zu ermöglichen. Die Software soll andere Teilnehmer im Netzwerk automatisch erkennen. Die Nachrichten werden mit der symmetrischen AES-Verschlüsselung verschlüsselt und ausgetauscht. Das Nachrichtenprotokoll baue ich von Grund auf selbst auf. Alle Änderungen werden mit GitHub protokolliert und lassen sich nachverfolgen.
\section{Methode}
In diesem Kapitel beschreibe ich die Methoden, die ich benutzt habe. Zuerst habe ich gründlich zu Netzwerkkommunikation, Verschlüsselung und der Programmiersprache Python recherchiert. Programmiert habe ich in VSCode zusammen mit uv.\citep{AstralshUv2026} Den Theorieteil habe ich mit \LaTeX{} geschrieben und die Quellen mit Zotero und BibTeX verwaltet.\citep{BetterBibTeXZotero} Alle Dateien sind in meinem GitHub-Repository abgespeichert: \url{https://github.com/Luka-Adjancic/Luka-Adjancic-Mauturaarbeit}.